\subsubsection{Contraste de la ventana dodecafásica con el centro recíproco}
\label{sec:contraste-ventana-reciproco}

La ventana finita \(\alpha_{12}^{(0)}\) de
\eqref{eq:alpha-ventana-cero}, obtenida con \(c_{13}=0\), admite además
el contraste histórico con el recíproco del centro impreso de la
constante inversa. Definamos
\[
 R_{2018}:=137.035999084,\qquad
 \alpha_{\mathrm{rec},2018}:=R_{2018}^{-1}.
\]
En esta operación se evalúa el recíproco del racional decimal
\(R_{2018}\). Esta referencia se distingue del centro tabulado
\(0.0072973525693\) utilizado en la tabla anterior.

\begin{center}
\begin{tabular}{@{}ll@{}}
\toprule
Expresión & Evaluación aritmética\\
\midrule
\(\alpha_{12}^{(0)}\) &
\(0.007297352569283800997285105472380663\)\\
\(\alpha_{\mathrm{rec},2018}\) &
\(0.007297352569283800997285105472380663567972560701\ldots\)\\
\(\alpha_{12}^{(0)}-\alpha_{\mathrm{rec},2018}\) &
\(\simeq-5.679725607012965\times10^{-37}\)\\
\((\alpha_{12}^{(0)})^{-1}-R_{2018}\) &
\(\simeq+1.066588006664435\times10^{-32}\)\\
\(\bigl((\alpha_{12}^{(0)})^{-1}-R_{2018}\bigr)/R_{2018}\) &
\(\simeq+7.783268730800000\times10^{-35}\)\\
\bottomrule
\end{tabular}
\end{center}

La relación decimal precisa es
\[
 \alpha_{12}^{(0)}
 =10^{-36}\left\lfloor10^{36}\alpha_{\mathrm{rec},2018}\right\rfloor.
\]
Por tanto, la coincidencia histórica corresponde al truncamiento a
treinta y seis posiciones decimales del recíproco; su redondeo a esa
posición termina en \(664\), mientras la ventana termina en \(663\).
Las diferencias de la tabla cuantifican el contraste entre ambos
números. La sección infinita \(\alpha_A\) conserva su normalización
compatible y su frontera propias, distinguidas de esta ventana finita.

El ajuste expresa la constante inversa como
\(137.035999084(21)\), con incertidumbre estándar
\(u(R_{2018})=2.1\times10^{-8}\)\cite{codata2018}.
Los decimales obtenidos al invertir su centro impreso pertenecen a la
evaluación aritmética de ese racional; la incertidumbre experimental
conserva la indicada por el ajuste. La comparación precedente con el
centro tabulado de \(\alpha\) mantiene su referencia y su residuo
normalizado propios.
